%
\section{Materials and Methods}
        \label{sec:methods}
%
\subsection{Experimental}

\subsubsection{Subjects}

\subsubsection{Apparatus and task}

\giampi{Placeholder: here Sur with Mattia (or someone from the hospital ) should provide technical details on the recordings.}
\subsubsection{Extraction and processing of neuronal data}
\giampi{Placeholder: here Sur with Mattia (or someone from the hospital ) should provide technical details on the recordings.}
\subsubsection{Behavioral analysis}
\giampi{Placeholder: Sur should provide details on the task and bhv metrics}

%
\subsection{Functional connectivity estimation}
    \label{ssec:methods_fc}
%
For each patient and recording phase \(\Phase \in \{\txtacr{rspre}, \txtacr{taskl}, \txtacr{taskt}, \txtacr{rspost}\}\) we estimate one weighted \gls{fc} network per canonical \gls{eeg} band. A frequency band is the interval \(\Band = [f_{\min}, f_{\max}]\) with \(0 \le f_{\min} < f_{\max} \le f_{\rm Ny}\), \(f_{\rm Ny} = f_s/2\); the six canonical bands (\(\delta\), \(\theta\), \(\alpha\), \(\beta\), \(\gamma_{\rm l}\), \(\gamma_{\rm h}\), Table~\ref{tab:frequency-bands}) are fixed once and reused for every patient and phase. All subsequent steps are performed independently per band, so that each phase yields a collection of weighted networks indexed by the triplet (patient, phase, band). The \acrshort{taskl} phase does not enter the primary inter-phase comparison of Section~\ref{ssec:methods_compare}, which is anchored on the \(\txtacr{rspre} \to \txtacr{taskt} \to \txtacr{rspost}\) triangle; it supplies the seen-only reference of the encoding/inference decomposition of Section~\ref{ssec:methods_arc}.\par
%
\begin{table}[!h]
    \centering
    \begin{tabularx}{\textwidth}{ X X c }
        \hline\hline
        \textbf{Band} & \textbf{Label} & \textbf{Range (Hz)}\\
        \hline
        Delta & \texttt{delta} & \(\qty{0.53}{}\)--\(\qty{4}{\hertz}\)\\
        Theta & \texttt{theta} & \(\qty{4}{}\)--\(\qty{8}{\hertz}\)\\
        Alpha & \texttt{alpha} & \(\qty{8}{}\)--\(\qty{13}{\hertz}\)\\
        Beta & \texttt{beta} & \(\qty{13}{}\)--\(\qty{30}{\hertz}\)\\
        Low gamma & \texttt{low\_gamma} & \(\qty{30}{}\)--\(\qty{80}{\hertz}\)\\
        High gamma & \texttt{high\_gamma} & \(\qty{80}{}\)--\(\qty{300}{\hertz}\)\\
        \hline\hline
    \end{tabularx}
    \caption{Canonical frequency bands used for \gls{fc} estimation.}
        \label{tab:frequency-bands}
\end{table}
%
Cross-spectral densities \(S_{ij}(f)\) between contacts \(i\) and \(j\) are estimated from the multichannel \gls{seeg} signal \(\TimeSeries \in \mathbb{R}^{N \times T}\) with Welch's averaged-periodogram method~\cite{welch1967use}, using a Hann window with \qty{50}{\percent} overlap and window length \texttt{nperseg} matched to the per-patient sampling rate so as to preserve a fixed frequency resolution \(\Delta f = \qty{0.5}{\hertz}\) across the cohort: \texttt{nperseg}\(= 4096\) samples for patients recorded at \qty{2048}{\hertz} and \texttt{nperseg}\(= 2048\) samples for the single patient recorded at \qty{1024}{\hertz}. The magnitude-squared coherence at frequency \(f\) is
\begin{equation}
    \MSC_{ij}(f) = \frac{\abs{S_{ij}(f)}^{2}}{S_{ii}(f) S_{jj}(f)} \in [0,1],
    \label{eq:methods_msc}
\end{equation}
and the coherency is the complex normalization \(\Coh_{ij}(f) = S_{ij}(f) / \sqrt{S_{ii}(f) S_{jj}(f)}\), so that \(\MSC_{ij}(f) = \abs{\Coh_{ij}(f)}^{2}\). The \gls{msc} captures every phase-consistent component of \(S_{ij}(f)\), including those produced by the recording montage rather than by neural interactions. Two such mechanisms are unavoidable on intracranial recordings of the type used here: \emph{volume conduction}, which propagates the field of a single source to nearby contacts effectively instantaneously, and the \emph{common reference}, which injects an identical electrical signal into every contact sharing a reference shaft. Both contributions enter \(S_{ij}(f)\) at zero phase lag and are therefore purely real, inflating \(\MSC_{ij}\) systematically for nearby (and especially same-shaft) pairs without reflecting any genuine inter-regional coupling.

Because the artifact contributions are purely real, the imaginary part of coherency removes them exactly. We therefore adopt the imaginary coherence~\cite{nolte2004identifying},
\begin{equation}
    \ImCoh_{ij}(f) \equiv \Im[\Coh_{ij}(f)] = \frac{\Im(S_{ij}(f))}{\sqrt{S_{ii}(f) S_{jj}(f)}} \in [-1,1],
    \label{eq:methods_imcoh_signed}
\end{equation}
whose immunity to zero-lag coupling is mathematically exact at every frequency, regardless of source amplitude, electrode geometry, or band. A nonzero \(\abs{\ImCoh_{ij}(f)}\) therefore certifies the presence of a genuine time-lagged interaction between contacts \(i\) and \(j\). For the purposes of constructing an \gls{lrg}-compatible adjacency matrix we use the band-averaged magnitude
\begin{equation}
    \Adjacency_{ij}(\Band) = \frac{1}{\abs{\mathcal{K}_{\Band}}} \sum_{k \in \mathcal{K}_{\Band}} \abs{\ImCoh_{ij}(f_{k})} \in [0,1], \qquad \mathcal{K}_{\Band} = \{k : f_{k} \in \Band\},
    \label{eq:methods_imcoh_bandavg}
\end{equation}
hereafter denoted \(\abs{\ImCoh}\). The resulting matrix is symmetric, non-negative, and has zero diagonal, satisfying the input constraints of the combinatorial Laplacian \(\lapl = \degm - \Adjacency(\Band)\) (Section~\ref{ssec:methods_lrg}) without any rectification step. Taking the magnitude before band-averaging discards the sign of \(\Im[\Coh_{ij}(f)]\), which encodes phase-lag directionality but is not needed for the undirected community-detection framework adopted here; this convention follows~\citet{ewald2012estimating,bastos2015functional}.
%
\subsection{Laplacian renormalization group framework}
    \label{ssec:methods_lrg}
%
For each (patient, phase, band) triplet, the adjacency matrix \(\Adjacency(\Band)\) of \eqref{eq:methods_imcoh_bandavg} defines a weighted undirected graph \(\graph = (\vertexs, \edges, \Adjacency)\) on \(N\) contacts. The combinatorial Laplacian
\begin{equation}
    \lapl = \degm - \Adjacency(\Band), \qquad \degm_{ii} = \sum_{j} \Adjacency_{ij}(\Band),
    \label{eq:methods_laplacian}
\end{equation}
is symmetric and positive semidefinite, with eigenvalues \(0 = \lambda_{1} < \lambda_{2} \le \cdots \le \lambda_{N} = \lambda_{\max}\) and orthonormal eigenvectors \(\{\vec{\phi}_{\ell}\}_{\ell=1}^{N}\). We use the combinatorial form throughout (rather than the random-walk Laplacian \(\lapl_{\rm RW} = \degm^{-1}\lapl\) or the symmetric-normalized form \(\lapl_{\rm sym} = \degm^{-1/2} \lapl \degm^{-1/2}\)) so that the diffusion semigroup is generated by a symmetric self-adjoint operator with a real-orthogonal spectral decomposition; \(\lapl_{\rm sym}\) is also symmetric \gls{psd} but diagonalizes a different operator from the combinatorial propagator that the cophenetic and subspace probes downstream both build on.

The \gls{lrg} construction~\cite{villegas2023laplacian,villegas2025multiscale} reads multiscale network organization through the heat semigroup generated by \(\lapl\). The Laplacian density matrix and its associated propagator are
\begin{equation}
    \dm(\tau) = \frac{e^{-\tau\lapl}}{\Tr[e^{-\tau\lapl}]}, \qquad \propag(\tau) = e^{-\tau\lapl},
    \label{eq:methods_density_propagator}
\end{equation}
with diffusion time \(\tau > 0\) acting as a continuous resolution parameter: small \(\tau\) resolves fine local structure (high-frequency Laplacian modes), large \(\tau\) integrates over progressively longer paths and washes out network heterogeneity. The von Neumann entropy \(\ent(\tau) = -\Tr[\dm(\tau) \log \dm(\tau)]\) interpolates between \(\ent \to \log N\) at \(\tau \to 0\) (uniform spectrum) and \(\ent \to 0\) at \(\tau \to \infty\) (collapse onto the trivial mode), and its logarithmic derivative
\begin{equation}
    C(\tau) = -\dv{\ent(\tau)}{\log_{10}\tau}
    \label{eq:methods_susceptibility}
\end{equation}
is the entropic susceptibility. On a network with a discrete spectral gap \(C(\tau)\) exhibits multiple peaks at the gap-induced crossover scales; the dense weight-heterogeneous \(\abs{\ImCoh}\) substrate has a continuous spectrum and therefore no such family of mesoscale peaks, yet it retains a single well-defined susceptibility peak \(\tau^{\ast}\) (the maximum of \(C(\tau)\)) at an intermediate scale of order \(10/\lambda_{\max}\), the characteristic scale at which the propagator collapses onto the trivial mode and inter-contact structure is lost. The natural resolution window of the propagator is \([\tau_{\min}, \tau^{\ast}]\), from the finest scale \(\tau_{\min} = 1/\lambda_{\max}\) resolved by \(\propag(\tau)\) up to this collapse scale, with the Fiedler time \(1/\lambda_{\rm gap}\) an interior scale of the window rather than its coarse bound. The cross-phase trace of Section~\ref{ssec:methods_compare} is anchored at the finest scale \(\tau = \tau_{\min} = 1/\lambda_{\max}\): the standard partition selector \(\Psi(n;\tau) = \mathcal{N}[\log_{10} \Delta_{n}(\tau) - \log_{10}\Delta_{n+1}(\tau)]\) of~\cite{villegas2025multiscale} isolates no privileged interior cut on this continuous-spectrum substrate, the linkage hierarchy at \(\tau_{\min}\) already folds every coarsening scale of \(\propag(\tau_{\min})\) into one per-pair image, and coarsening a \emph{cross-phase} comparison toward \(\tau^{\ast}\) drives the diffusion geometry into the collapse regime, where a phase-invariant baseline dominates every phase alike. A within-phase read-out is not subject to this last constraint and is read instead at the coarse end of the window (Section~\ref{ssec:methods_epi}).

Two coupled multiscale objects are extracted from the propagator \(\propag(\tau_{\min})\) at the working scale and serve, in the rest of the analysis, as the probes of the network's information-communication geometry. The first is the \emph{cophenetic communication distance} \(\Dcoph\) between contact pairs, defined as follows. The density matrix \(\dm(\tau_{\min})\) of \eqref{eq:methods_density_propagator} induces a per-pair communication distance
\begin{equation}
    \distm_{ij}(\tau_{\min}) = \frac{1 - \delta_{ij}}{\dm_{ij}(\tau_{\min})}
    \label{eq:methods_distm}
\end{equation}
at the working scale~\cite{villegas2025multiscale}. Average-linkage agglomerative clustering~\cite{kaufman2009finding} on \(\distm(\tau_{\min})\) --- chosen as the finite-sample estimator of the mean cophenetic distance in the continuous-spectrum regime, since single-linkage produces chained dendrograms unstable to noise in the heavy-tailed \(\distm(\tau_{\min})\) distribution and Ward minimizes a within-cluster variance under a Gaussian assumption that does not match the \(1/\dm_{ij}\) distance scale --- produces a deterministic linkage hierarchy \(\mathcal{L}(\tau_{\min})\) whose \(N-1\) merge heights \(\{h_{n}\}_{n=1}^{N-1}\) span the coarsening scales of the network as resolved by \(\propag(\tau_{\min})\). The cophenetic distance reads, for each leaf pair \((i,j)\),
\begin{equation}
    \Dcoph_{ij} = h\bigl(\mathrm{LCA}_{\mathcal{L}(\tau_{\min})}(i,j)\bigr),
    \label{eq:methods_distm_coph}
\end{equation}
the merge height at which \(i\) and \(j\) first coalesce in \(\mathcal{L}(\tau_{\min})\), where \(\mathrm{LCA}_{\mathcal{L}(\tau_{\min})}(i,j)\) is the lowest common ancestor of \(i\) and \(j\) in the linkage. \(\Dcoph\) is ultrametric by construction (the cophenetic image of any agglomerative linkage is ultrametric) and reads, pair by pair, the scale at which the two contacts coalesce into a single communication unit. This is the per-pair image of the multiscale structure of \(\propag(\tau_{\min})\): the dimensionality remains \(N(N-1)/2\) pair values, but each value is now evaluated at the intrinsic communication-merge scale of that pair rather than at a single fixed \(\tau\). The raw distance \(\distm(\tau_{\min})\) is not read directly: for our continuous Laplacian spectrum the linkage hierarchy replaces the spectral-gap scale-identification mechanism of~\cite{villegas2025multiscale}, and the cophenetic image \(\Dcoph\) is the natural multiscale summary of \(\distm\) along the merge-height continuum. The second object is the \emph{leading-eigenmode subspace} \(U_{k} \in \mathbb{R}^{N \times k}\) whose columns are the eigenvectors associated with the \(k\) smallest non-zero eigenvalues of \(\lapl\), \(U_{k} = [\vec{\phi}_{2}  \vec{\phi}_{3}  \cdots  \vec{\phi}_{k+1}]\). These are the slowest non-trivial diffusion modes and span the subspace in which low-frequency communication is concentrated.

The two objects read structurally distinct aspects of the same information-communication geometry: \(\Dcoph\) reads the per-pair hierarchical-merge structure (which contacts coalesce at which scale), and \(U_{k}\) reads the subspace alignment of the slow diffusion modes at a fixed cutoff \(k\). They are not derivable from each other --- the linkage hierarchy is not recoverable from \(U_{k}\) alone, and the eigenmode subspace is not recoverable from \(\Dcoph\) alone --- and cross-phase comparison at the two levels reads two complementary aspects of any coordinated reorganization (Section~\ref{ssec:methods_compare}).
%
\subsection{Comparison of the \glsentryshort{lrg} information-communication output proxies}
    \label{ssec:methods_compare}
%
The trace is probed two ways on the graph, and read against the untransformed substrate. The primary read-out is the network's \gls{lrg} multiscale communication geometry, taken per contact pair as the cophenetic distance \(\rhocoph\) (Section~\ref{sssec:methods_compare_ctm}). Against it we set a conventional spectral counterpart --- the rotation of the leading Laplacian eigenmode subspace, the representation on which standard spectral-clustering pipelines rely, measured by the Grassmann distance \(d_{\rm G}\) (Section~\ref{sssec:methods_compare_grassmann}). Both act on the \(\abs{\ImCoh}\) network \(\Adjacency(\Band)\), whose own per-pair overlap \(\rhoraw\) (Section~\ref{sssec:methods_compare_rawfc}) is the pre-transform reference against which the multiscale step is measured. Each read-out yields a per-(patient, band) trace scalar sharing one sign orientation --- positive is the \emph{trace direction} (an \acrshort{rspost} configuration retaining a residual imprint of the \acrshort{taskt} configuration), negative the \emph{reset direction} (an \acrshort{rspost} returned to baseline rest) --- and each removes the patient's static, phase-invariant connectivity backbone in the manner native to its object.
%
The two per-pair read-outs, \(\rhocoph\) and \(\rhoraw\), act on per-pair vectors that carry the static backbone additively and remove it by explicit subtraction: a task-induced and a rest-induced per-pair shift vector, each referenced to an independent \acrshort{rspre} half, are compared through their Spearman correlation. Being rank-based, this correlation reads the \emph{direction} of the per-pair reorganization --- whether task and rest displace the same pairs in the same rank order --- and is blind to its overall magnitude. \(\rhocoph\) applies this to the \gls{lrg} multiscale cophenetic distance and is the primary read-out on which the trace rests; \(\rhoraw\) applies the identical statistic to the raw edge weights, so that the difference between them isolates what the multiscale transform contributes over the untransformed substrate.
%
The Grassmann read-out instead compares whole subspaces: the leading Laplacian eigenmodes are swept across the full spectral range, and \(d_{\rm G}\) measures the \acrshort{rspre}-to-\acrshort{taskt} rotation of that subspace against the \acrshort{taskt}-to-\acrshort{rspost} rotation, so baseline enters as a reference scale rather than by subtraction. Being a distance, it reads the \emph{magnitude} of whole-network subspace rotation rather than a pair-resolved pattern. The cophenetic and Grassmann read-outs are built from the same Laplacian eigendecomposition and are therefore correlated, not independent, but they are different reductions of it: the cophenetic integrates all merge-height scales of the diffusion propagator into one nested per-pair hierarchy, while the Grassmann keeps only the flat span of the leading modes that a conventional spectral embedding retains. Where the two agree, the agreement certifies that the trace is robust to how the spectrum is read; where they dissociate --- a multiscale per-pair trace with no leading-mode rotation, or a leading-mode rotation with no persistent per-pair hierarchy --- the \gls{lrg} multiscale geometry and the conventional spectral summary are resolving different reorganizations, each seeing structure the other misses.
%
\subsubsection{Per-pair split-baseline trace probes}
    \label{sssec:methods_compare_perpair}
%
The two per-pair comparisons share a single split-baseline scheme and differ only in the per-pair object they act on. We define the scheme once on a generic symmetric per-pair object \(O\), with \(O^{\Phase}\) its value at phase \(\Phase\) and \(O^{\Phase}_{ij}\) its \(N(N-1)/2\) upper-triangular entries indexing contact pairs \((i,j)\); the instantiations of Sections~\ref{sssec:methods_compare_rawfc} and~\ref{sssec:methods_compare_ctm} set \(O\) to the raw \(\abs{\ImCoh}\) adjacency and to the \gls{lrg} cophenetic distance respectively.

For each (patient, band) cell, \acrshort{rspre} is split into two halves \(A\) and \(B\) of equal duration and \(O\) is computed independently on each, giving two independent baselines \(O^{\txtacr{rspre}_{A}}\) and \(O^{\txtacr{rspre}_{B}}\). We form the task-side and rest-side per-pair shifts against \emph{different} halves,
\begin{equation}
    \Delta_{\rm task}^{A}(i,j) = O^{\txtacr{taskt}}_{ij} - O^{\txtacr{rspre}_{A}}_{ij}, \qquad \Delta_{\rm rest}^{B}(i,j) = O^{\txtacr{rspost}}_{ij} - O^{\txtacr{rspre}_{B}}_{ij},
    \label{eq:methods_perpair_shifts}
\end{equation}
so that the shared-baseline term --- which would correlate the two shift maps trivially if both subtracted the same \acrshort{rspre} estimate --- is removed by construction. The single-arm trace scalar is their Spearman correlation across all pairs,
\begin{equation}
    \rho_{\mathrm{split}} = \rho_{S}\!\left(\Delta_{\rm task}^{A},\, \Delta_{\rm rest}^{B}\right),
    \label{eq:methods_perpair_rho}
\end{equation}
with \(\rho_{\mathrm{split}} > 0\) the trace direction: the pairs the task displaces are, in rank, the pairs displaced in post-task rest.

The assignment ``\(A\!\to\!\)task-side, \(B\!\to\!\)rest-side'' in \eqref{eq:methods_perpair_rho} is, however, arbitrary; the mirror assignment is equally valid and defines a second, equally-legitimate arm \(\rho_{\mathrm{split}}' = \rho_{S}(O^{\txtacr{taskt}} - O^{\txtacr{rspre}_{B}},\, O^{\txtacr{rspost}} - O^{\txtacr{rspre}_{A}})\). Nothing prefers one arm over the other, yet the two can disagree: because the halves differ on roughly \qty{45}{\percent} of pairs, low-reorganization (near-zero) patients can receive \emph{opposite signs} from the two arms, so a per-patient trace/reset label would depend on an arbitrary bookkeeping choice. The estimator of record therefore \emph{symmetrizes} the two arms, averaging \(\rho_{\mathrm{split}}\) with its mirror \(\rho_{\mathrm{split}}'\); the average is invariant to the \(A\!\leftrightarrow\!B\) half-labelling by construction, uses the same data and the same matched-strength surrogate (Section~\ref{sssec:methods_compare_stats}), and leaves strong tracers essentially unchanged while removing the arbitrary-half artifact for the near-zero patients. Instantiated on the cophenetic distance and on the raw substrate, this arm-symmetric estimator defines the primary trace \(\rhocoph\) (Section~\ref{sssec:methods_compare_ctm}) and the substrate reference \(\rhoraw\) (Section~\ref{sssec:methods_compare_rawfc}); those sections give its explicit symmetric form, and throughout \(\rhocoph\) and \(\rhoraw\) denote the symmetrized estimators. Per patient we report the arm-symmetric estimate together with the half-arm disagreement \(\tfrac{1}{2}\abs{\rho_{\mathrm{split}} - \rho_{\mathrm{split}}'}\) as its standard error, labelling any patient whose estimate is smaller than one such SE \emph{undetermined} --- its sign is estimator noise rather than a biological verdict. The single-arm \(\rho_{\mathrm{split}}\) is retained only as a supplementary column.

Collapsing the two halves into a single shared \acrshort{rspre} baseline would remove the half-choice outright, but was rejected because it reintroduces exactly the correlated-error inflation the split design exists to prevent; the arm-symmetric estimator is thus the free refinement of the split design, not a retreat from it. The rank-based Spearman form is used throughout because the trace question is one of pattern, not magnitude --- whether task and rest reorganize the same pairs in the same order --- so the statistic is deliberately orthogonal to the edge-strength backbone and robust to the heavy-tailed per-pair shift distribution; the object-specific motivation on \(\Dcoph\) is given at its instantiation.

The cohort verdict rests on one mandatory null and two supplementary controls (cohort tests in Section~\ref{sssec:methods_compare_stats}), stated here for the cophenetic trace \(\rhocoph\); the substrate \(\rhoraw\) is gated identically. The \textbf{gate} is the matched-strength surrogate null: each phase matrix is rewired under the strength-preserving \(4\)-cycle \(\pm\delta\) scheme (\(R = 200\) surrogates; Section~\ref{sssec:methods_compare_stats}), \(\rhocoph\) is recomputed on every surrogate to build a per-(patient, band) null distribution, and a \(\rhocoph\) exceeding its patient-matched surrogate median is not explained by strength alone. The first \textbf{control} is the drift-floor null, computed without surrogates: splitting \acrshort{rspost} into two equal-duration halves as well, the rest-only correlation
\begin{equation}
    \rho_{\rm drift} = \rho_{S}\!\left(O^{\txtacr{rspre}_{B}} - O^{\txtacr{rspre}_{A}},\; O^{\txtacr{rspost}_{B}} - O^{\txtacr{rspost}_{A}}\right)
    \label{eq:methods_perpair_drift}
\end{equation}
is the task-free counterpart of \(\rhocoph\), and the trace is required to exceed its own drift floor so that a monotone within-session drift cannot fake it. The second \textbf{control} is the cross-probe restriction \(\rho_{\rm xprobe}\), which recomputes the trace on cross-probe pairs alone (pairs whose two contacts sit on different electrode shafts); the trace must not degrade significantly under this restriction, checking that it is not carried by same-probe short-range coupling --- already suppressed under \(\abs{\ImCoh}\), so this is a robustness control rather than a de-biasing requirement.
%
\subsubsection{Substrate-level instantiation on \texorpdfstring{\(\Adjacency(\Band)\)}{W(B)}}
    \label{sssec:methods_compare_rawfc}
%
The substrate-level instantiation sets \(O = \Adjacency(\Band)\), the raw \(\abs{\ImCoh}\) adjacency before any \gls{lrg} transformation, so that the arm-symmetric estimator of Section~\ref{sssec:methods_compare_perpair} evaluated on the untransformed edge weights defines the substrate overlap \(\rhoraw \equiv \rhoraw_{\mathrm{sym}}\),
\begin{equation}
    \rhoraw = \tfrac{1}{2}\Big[\rho_{S}\!\left(\Adjacency^{\txtacr{taskt}} - \Adjacency^{\txtacr{rspre}_{A}},\; \Adjacency^{\txtacr{rspost}} - \Adjacency^{\txtacr{rspre}_{B}}\right) + \big(A\!\leftrightarrow\!B\big)\Big],
    \label{eq:methods_rho_raw}
\end{equation}
where \((A\!\leftrightarrow\!B)\) denotes the mirror arm with the two \acrshort{rspre} halves exchanged, and the drift floor and cross-probe restriction are defined by the same substitution. This substrate overlap treats each upper-triangular entry of \(\Adjacency(\Band)\) as an independent observation and asks whether the same contact pairs are coupled across phases at the level of direct \(\abs{\ImCoh}\) edge weights. It is the pre-transform reference against which the multiscale step is measured: comparing \(\rhocoph\) against \(\rhoraw\) isolates what the \gls{lrg} multiscale transform adds over the untransformed substrate.
%
\subsubsection{Cophenetic-level instantiation on \texorpdfstring{\(\Dcoph\)}{D coph}}
    \label{sssec:methods_compare_ctm}
%
The cophenetic instantiation is the primary, multiscale read-out: the arm-symmetric estimator of Section~\ref{sssec:methods_compare_perpair} evaluated on \(O = \Dcoph\), the cophenetic communication distance of \eqref{eq:methods_distm_coph}, defines the cophenetic trace of record, \(\rhocoph \equiv \rhocoph_{\mathrm{sym}}\), which the main text uses throughout,
\begin{equation}
    \rhocoph = \tfrac{1}{2}\Big[\rho_{S}\!\left(\Dcoph^{\txtacr{taskt}} - \Dcoph^{\txtacr{rspre}_{A}},\; \Dcoph^{\txtacr{rspost}} - \Dcoph^{\txtacr{rspre}_{B}}\right) + \big(A\!\leftrightarrow\!B\big)\Big],
    \label{eq:methods_rho_coph}
\end{equation}
its single-arm value \(\rhocoph_{\mathrm{split}}\) retained only as a supplement, and the drift floor and cross-probe restriction defined by the same substitution under the matched-strength gate of Section~\ref{sssec:methods_compare_stats}. Each per-pair cophenetic image is rebuilt by rerunning the full \gls{lrg} pipeline of Section~\ref{ssec:methods_lrg} --- eigendecomposition, propagator at \(\tau_{\min} = 1/\lambda_{\max}\), average-linkage agglomeration, cophenetic image --- independently on each \acrshort{rspre} and \acrshort{rspost} half, so that \(\Dcoph^{\txtacr{rspre}_{A}}\) and \(\Dcoph^{\txtacr{rspre}_{B}}\) are independent cophenetic images on disjoint samples. On \(\Dcoph\) the Spearman form is additionally motivated by the tie structure: the cophenetic image takes only \(N-1\) distinct values \(\{h_{n}\}_{n=1}^{N-1}\), so its informative content is the rank ordering of merge heights rather than their precise magnitudes, and rank statistics handle the heavy ties cleanly and are invariant to the monotone rescaling that the cophenetic map already imposes.

The cophenetic image \(\Dcoph\) is read in place of the fixed-scale distance \(\distm(\tau_{\min})\). The latter reads pairwise dissimilarity at the single scale \(\tau_{\min}\); \(\Dcoph\) instead reads, pair by pair, the merge scale at which two contacts coalesce into one communication unit, so the per-pair correlation on \(\Dcoph\) captures cross-phase persistence of the \emph{multiscale-resolved} organization rather than single-scale persistence only. As established in Section~\ref{ssec:methods_lrg}, the continuous Laplacian spectrum of the \(\abs{\ImCoh}\) substrate precludes the gap-based \(\tau\)-sweep reading of~\cite{villegas2025multiscale}, and the linkage hierarchy \(\mathcal{L}(\tau_{\min})\) is the multiscale carrier that replaces it.
%
\subsubsection{Subspace distance on the leading eigenmodes}
    \label{sssec:methods_compare_grassmann}
%
The leading-mode subspace \(U^{\Phase}_{k} = \mathrm{span}\{\vec{\phi}_{2}, \ldots, \vec{\phi}_{k+1}\}\) defined in Section~\ref{ssec:methods_lrg} is a point on the Grassmann manifold \(\mathrm{Gr}(k, N)\), and two phases are compared through the chordal Grassmann distance. For orthonormal column matrices \(A, B \in \mathbb{R}^{N \times k}\) representing the two subspaces,
\begin{equation}
    d_{\rm G}(A, B)^{2} = k - \norm{A^{\transp} B}_{F}^{2} = \sum_{i=1}^{k} \sin^{2}\theta_{i},
    \label{eq:methods_grassmann}
\end{equation}
where \(\theta_{1} \le \cdots \le \theta_{k}\) are the principal angles between the two subspaces; the chordal form is gauge-invariant under sign flips of any eigenvector and under basis rotations within degenerate eigenspaces, and the trace identity on the right makes the full \(k\)-sweep computable from a single cross-correlation \(A^{\transp} B\). We sweep the cutoff over \(k \in \{2, \ldots, 112\}\) --- the entire non-trivial spectral range, not a handful of leading modes --- so that a reorganization at any resolution scale can register (the upper bound matches \(N - 1\) of the smallest-\(N\) patient). The per-patient, per-\(k\) trace statistic is the phase-triangle
\begin{equation}
    T_{\rm G}(k; s, b) = d_{\rm G}\!\left(U^{\txtacr{rspre}_{A}}_{k}, U^{\txtacr{taskt}}_{k}\right) - d_{\rm G}\!\left(U^{\txtacr{taskt}}_{k}, U^{\txtacr{rspost}}_{k}\right),
    \label{eq:methods_T_G}
\end{equation}
where \(s\) indexes patients and the early \acrshort{rspre} half is used; \(T_{\rm G}(k; s, b) > 0\) is the trace direction at cutoff \(k\), in the shared sign orientation of Section~\ref{ssec:methods_compare}.

The cutoff \(k\) is a nuisance dimension: a trace can land at any resolution scale, and reading off any individual \(k\) would both multiply-test order-\(100\) hypotheses and discard the across-\(k\) structure that a strength-preserving null cannot fake. We therefore collapse the \(k\)-axis with a single cluster-extent permutation test. At each \(k\), a one-sided paired Wilcoxon compares the observed \(T_{\rm G}(k)\) against the patient-matched matched-strength surrogates of Section~\ref{sssec:methods_compare_stats}, yielding \(p_{k}(b)\); the band-level statistic is the cluster mass \(T_{\rm G}^{\ast}(b) = \sum_{k : p_{k}(b) < 0.05}\left(-\log_{10} p_{k}(b)\right)\), which sums evidence over \emph{all} significant cutoffs and so responds to both across-\(k\) extent and per-\(k\) depth, robust to fragmentation of a contiguous run without an auxiliary contiguity threshold. Its significance \(\pmassemp(b)\) is the empirical \(p\)-value of \(T_{\rm G}^{\ast}(b)\) against a phantom-surrogate permutation null built from the matched-strength ensemble, and the band clears the Grassmann gate when \(\pmassemp(b) < 0.05\) and the verdict is leave-one-patient-out robust. The \([0,1]\) normalization of \(T_{\rm G}^{\ast}\), a per-patient analogue reported descriptively, and the phantom-null construction are given in the Supplement; \(T_{\rm G}^{\ast}\) is a significance-weighted extent statistic, not a trace amplitude, so a band that fails the gate carries no interpretable magnitude.

As set out in Section~\ref{ssec:methods_compare}, \(\rhocoph\) and \(d_{\rm G}\) are built from the same Laplacian eigendecomposition --- the cophenetic from the multiscale diffusion propagator, the Grassmann from the leading-mode subspace of a conventional spectral embedding --- so a band clearing both is read as robust to how the spectrum is summarized, and a band-by-band dissociation as the two summaries resolving different reorganizations. No combined cross-probe scalar is formed.
%
\subsubsection{Statistical inference and robustness}
    \label{sssec:methods_compare_stats}
%
\paragraph{Primary gate --- matched-strength surrogate null.}
Every trace verdict rests on one mandatory null: a strength-preserving surrogate that removes the per-node coupling backbone while randomizing everything else. For each (patient, band, phase) cell we generate \(R = 200\) surrogates of the adjacency \(\Adjacency(\Band)\) by repeated \(4\)-cycle \(\pm\delta\) edge-weight rewirings, which preserve the per-node strength sequence \(s_{i} = \sum_{j} \Adjacency_{ij}(\Band)\) to machine precision and keep every weight within the original \([0, 1]\) interval while the marginal edge-weight distribution drifts. The probe of interest is recomputed on each surrogate --- rerunning the full \gls{lrg} pipeline of Section~\ref{ssec:methods_lrg} for \(\rhocoph\) and the Grassmann distance, or truncated at the substrate level for \(\rhoraw\) --- to yield a per-(patient, band) null at matched strength, with the symmetric, arbitrary-half-free construction of \(\rhocoph\) applied identically to observation and surrogate. Cohort significance for the cophenetic trace is a one-sided paired Wilcoxon of the observed per-patient \(\rhocoph\) against its patient-matched surrogate median, under \(H_{1}\): observed in the trace direction; the same test on the substrate reference \(\rhoraw\) provides the pre-transform comparison. Because the rewiring strips the strength structure that inflates any correlation of edge levels, a \(\rhocoph\) surviving this gate is not explained by strength alone. This is the minimum null on which any cohort claim rests.

\paragraph{Supplementary per-pair controls.}
Two surrogate-free controls, read at the cophenetic layer alongside the gate, are defined in Section~\ref{sssec:methods_compare_perpair}. The \emph{drift-floor null} \(\rhodrift\) of \eqref{eq:methods_perpair_drift} is the task-free counterpart of the trace, and the cohort test is a one-sided paired Wilcoxon under \(H_{1}: \rhocoph > \rhodrift\), requiring the observed correlation to exceed its own within-session drift. The \emph{cross-probe restriction} \(\rhoxprobe\) recomputes the trace on cross-shaft pairs only, and the cohort test is a non-degradation gate: the one-sided paired Wilcoxon on \(\rhocoph - \rhoxprobe\) must be non-significant, with the cohort-median sign of \(\rhoxprobe\) matching that of \(\rhocoph\).

\paragraph{Grassmann band-level gate.}
The Grassmann probe is gated at the band level by the cluster-extent permutation test of Section~\ref{sssec:methods_compare_grassmann} --- \(\pmassemp(b) < 0.05\), leave-one-out robust --- which collapses the \(k\)-sweep and is itself the family-level correction over \(k\), so no separate within-family correction is applied at that layer. Full construction of the cluster-mass statistic and its phantom-surrogate null is deferred to the Supplement.

\paragraph{Robustness diagnostics.}
Every cohort-level gate is paired with a leave-one-out (\gls{loo}) sensitivity diagnostic: the test is recomputed \(n = 10\) times, each excluding one patient, and the maximum resulting \(p\)-value (\gls{loo} max \(p\)) is reported alongside the cohort \(p\) and the patient whose exclusion drives it. \gls{loo} max \(p\) is descriptive and never gates a verdict; verdicts whose \gls{loo} max \(p\) crosses \num{0.05} are flagged in the Results as single-patient-leveraged. A single \emph{epileptogenic-zone exclusion} regime drops each patient's clinically identified epileptogenic contacts~\cite{cardinale2013stereoelectroencephalography} and recomputes each probe on the residual set --- paired with a node-count decimation control, since rebuilding the \gls{lrg} hierarchy on any reduced contact set shifts \(\rhocoph\) --- and is reported for the cophenet-positive bands, with the Grassmann gate re-run on the epi-excluded surrogate ensemble where applicable. Finally, the between-patient spread in \(\rhocoph\) is checked against measurement reliability: each phase's cophenetic image is rebuilt on two disjoint halves of the recording and the split-half Spearman of the two taken as the reliability of the estimate, against which the between-patient spread in trace strength is referred.

The cohort verdict of both read-outs across all six bands is collected in Table~\ref{tab:probe_band_full}.

\begin{table}[t]
    \centering
    \caption{\textbf{Full six-band probe map.} Cohort verdicts on the cophenetic
      (\(\rhocoph\); multiscale, nested across scales) and Grassmann (\(d_{G}\); leading
      global modes at a single scale) read-outs for all six bands, each against its own
      strength-preserving matched-strength control (\(p\) from the paired cohort test).
      \(\checkmark\) marks a cohort-level trace that clears that read-out's gate; --- marks
      none. \(\beta\) is the only band both read-outs register; \(\alpha\) is
      cophenetic-only and \(\lowgamma\) Grassmann-only, so each probe recovers a
      reorganization the other misses. The four unambiguous bands are shown in the main
      text (Table~\ref{tab:probe_band}).}
    \label{tab:probe_band_full}
    \begin{tabular}{lcc}
        \toprule
        Band & Cophenetic \(\rhocoph\) & Grassmann \(d_{G}\) \\
        \midrule
        \(\beta\)      & \(\checkmark\)~~(\(p = 0.032\)) & \(\checkmark\)~~(\(p = 0.005\)) \\
        \(\alpha\)     & \(\checkmark\)~~(\(p = 0.024\)) & ---~~(\(p = 0.35\)) \\
        \(\lowgamma\)  & ---~~(\(p = 0.08\))            & \(\checkmark\)~~(\(p = 0.005\)) \\
        \(\delta\)     & ---~~(\(p = 0.08\))            & (\(\checkmark\))\textsuperscript{$\dagger$}~~(\(p = 0.005\)) \\
        \(\highgamma\) & ---~~(\(p = 0.08\))            & ---~~(\(p = 0.06\)) \\
        \(\theta\)     & ---~~(\(p = 0.78\))            & ---~~(\(p = 0.16\)) \\
        \bottomrule
    \end{tabular}

    \smallskip
    {\footnotesize \textsuperscript{$\dagger$}\(\delta\) clears the Grassmann cohort mass
    gate but is not leave-one-out robust (\gls{loo} max \(p = 0.055\)); a weak trace.}
\end{table}

\subsection{Anatomical localization of the trace}
\label{sssec:methods_anatomy}

The per-pair trace is localized to anatomy by decomposing it into per-contact contributions and asking which anatomical systems concentrate the trace above the patient's own brain-wide level. The same construction localizes either cross-phase trace, differing only in the task-side phase --- the test phase for the primary trace, or the learning phase for the encoding anchor of Section~\ref{ssec:methods_arc}. Writing the task-side and rest-side per-pair shifts \(\Delta_{\rm task}^{A}\) and \(\Delta_{\rm rest}^{B}\) of Section~\ref{sssec:methods_compare_perpair} in centered-rank form, the signed contribution of pair \((i,j)\) to the single-arm Spearman trace is
\begin{equation}
    c^{A}(i,j) = \Big[\rank(\Delta_{\rm task}^{A})_{ij} - \tfrac{P+1}{2}\Big]\Big[\rank(\Delta_{\rm rest}^{B})_{ij} - \tfrac{P+1}{2}\Big],
    \label{eq:methods_anat_concordance}
\end{equation}
ranked over all \(P = N(N-1)/2\) contact pairs and normalized so that \(\sum_{i<j} c^{A}(i,j)\big/\big[P(P^{2}-1)/12\big] = \rho_{\mathrm{split}}\) returns the single-arm cophenetic trace. The contribution of record is the arm-symmetric concordance, averaging the two \acrshort{rspre}-half assignments exactly as the trace itself (Section~\ref{sssec:methods_compare_ctm}),
\begin{equation}
    c(i,j) = \tfrac{1}{2}\Big[c^{A}(i,j) + c^{B}(i,j)\Big],
    \label{eq:methods_anat_sym}
\end{equation}
with \(c^{B}\) the mirror concordance with the two halves exchanged, so that \(\sum_{i<j} c(i,j)\big/\big[P(P^{2}-1)/12\big] = \rhocoph\) recovers the symmetric cophenetic trace of record and \(c^{A}\) is retained only as a supplement; every subsequent step is built on the symmetric \(c\). Each pair carries its concordance to both of its endpoint contacts, a contact's score is the mean concordance over the pairs incident to it, and this score is demeaned within the patient over its anatomically labeled contacts,
\begin{equation}
    \tilde{u}_{i} = u_{i} - \frac{1}{|V|}\sum_{k \in V} u_{k}, \qquad u_{i} = \frac{1}{N-1}\sum_{j \neq i} c(i,j),
    \label{eq:methods_anat_node}
\end{equation}
so that the statistic measures where the trace \emph{concentrates} relative to the patient's brain-wide average rather than merely where it is present --- a hotspot on a distributed trace, not a container. Each contact is assigned a \gls{dk} region from its implantation tissue labels (white-matter, subcortical, and unlabeled contacts are excluded from the anatomical universe \(V\)) and thence to one of nine a-priori cortical systems --- medial-temporal, lateral-temporal, orbitofrontal, prefrontal, sensorimotor, parietal, occipital, cingulate, and insula --- with \gls{ofc} pooling the medial and lateral orbitofrontal regions across hemispheres. The per-patient system mean and its cohort statistic are
\begin{equation}
    V_{g}^{(s)} = \frac{1}{|C_{g,s}|}\sum_{i \in C_{g,s}} \tilde{u}^{(s)}_{i}, \qquad M_{g} = \operatorname{median}_{\,s\,:\,C_{g,s} \neq \varnothing}\; V_{g}^{(s)},
    \label{eq:methods_anat_system}
\end{equation}
where \(C_{g,s}\) is the set of contacts of patient \(s\) in system \(g\) and the cohort median runs over the patients implanted there; in the shaft-collapsed variant each electrode shaft is first reduced to its mean incidence and \(V_{g}^{(s)}\) averages over distinct shafts, removing within-shaft pseudoreplication.

The per-system statistic is gated against the same strength-preserving matched-strength surrogate as the trace (\(4\)-cycle \(\pm\delta\) rewiring preserving every node's strength, \(R\) surrogates; Section~\ref{sssec:methods_compare_stats}), re-running the entire symmetric-concordance-to-cohort-median chain of \eqref{eq:methods_anat_concordance}--\eqref{eq:methods_anat_system} on each surrogate to build the null \(\{M_{g}^{r}\}_{r=1}^{R}\). The upper- and lower-tail empirical \(p\)-values
\begin{equation}
    p^{\rm up}_{g} = \frac{1 + \#\{r : M_{g}^{r} \ge M_{g}\}}{R+1}, \qquad p^{\rm lo}_{g} = \frac{1 + \#\{r : M_{g}^{r} \le M_{g}\}}{R+1}
    \label{eq:methods_anat_pval}
\end{equation}
identify systems that carry the trace (upper tail) and systems significantly depleted of it (lower tail); because the surrogate preserves each node's strength exactly, a surviving system is a topological concentration of the trace and not a consequence of hub geometry. Within the band under test the two tails are separately \gls{bh}--\gls{fdr}-corrected across the covered a-priori systems --- the multi-system family being the coordinated unit of inference --- and a system localizes the trace when its corrected \(q < \num{0.05}\).

\paragraph{Localization robustness.}
The localization is reported under four conditions --- the factorial of epileptogenic-zone \{inclusion, exclusion\} with \{contact-level, shaft-collapsed\} aggregation --- epi-inclusion being primary and the remaining three corroborative. Here epileptogenic-zone exclusion removes the seizure-onset contacts from the anatomical universe \(V\) --- the demean of \eqref{eq:methods_anat_node} and the per-system aggregation --- while the symmetric concordance itself is computed on the full graph in both conditions; no diffusion hierarchy is rebuilt on a reduced contact set, so the node-count decimation control of Section~\ref{sssec:methods_compare_stats} does not apply at this layer. The localization is required to be leave-one-patient-out robust: the cohort median and its surrogate \(q\) are recomputed dropping each implanted patient in turn, and the verdict stands only when the worst-case \gls{loo} \(q\) remains below \num{0.05}. A within-patient permutation of the region labels provides a coarse atlas check but does not gate the verdict; the matched-strength surrogate of \eqref{eq:methods_anat_pval} is the gate, as everywhere in this work.

\subsection{Tissue stratification of the trace}
\label{ssec:methods_tissue}

These recordings come from epilepsy patients, implanted because of pathology, so a persistence attributed to reasoning could instead ride the epileptogenic network or the white-matter contacts that motivated the implant. Attributing the trace to tissue means asking not whether it is present but which tissue carries it. We classify every contact by two independent binary labels read from the implantation record --- a structural label (cortical gray matter versus white matter) from its \gls{dk} tissue assignment, and a pathological label (inside versus outside the clinically defined seizure-onset zone \(E\)) --- and sort every contact pair by the labels of its two endpoints. Each partition yields two within-class cells and a crossing cell: gray--gray, white--white, and gray--white for the structural partition; \gls{soz}-internal (both endpoints in \(E\)), extra-\gls{soz} (neither), and \gls{soz}-crossing for the pathological partition.

For a pair class \(\mathcal{T}\) --- the subset of contact pairs whose endpoints satisfy the class rule --- the within-class trace is the arm-symmetric cophenetic estimator of Sections~\ref{sssec:methods_compare_perpair} and~\ref{sssec:methods_compare_ctm} evaluated on the pairs of that class alone,
\begin{equation}
    \rhocoph_{\mathcal{T}} = \tfrac{1}{2}\Big[\rho_{S}\!\left(\Delta_{\rm task}^{A}\big|_{\mathcal{T}},\; \Delta_{\rm rest}^{B}\big|_{\mathcal{T}}\right) + \big(A\!\leftrightarrow\!B\big)\Big],
    \label{eq:methods_tissue_class}
\end{equation}
where \(\,\cdot\,\big|_{\mathcal{T}}\) restricts a per-pair shift to the pairs of \(\mathcal{T}\) and \((A\!\leftrightarrow\!B)\) is the mirror arm of Section~\ref{sssec:methods_compare_ctm}. The cophenetic images \(\Dcoph^{\Phase}\) are built once on the full contact set (Section~\ref{ssec:methods_lrg}); only the per-pair shift vectors are subset to the class, so no diffusion hierarchy is rebuilt on a tissue subset and \(\rhocoph_{\mathcal{T}}\) is the whole-network trace read through the pairs of one tissue class, with positive values in the trace direction.

Two questions gate a class: whether it carries more trace than an equal-size set of pairs anywhere in the network, and whether that concentration survives node strength. The first is a pair-count control. Drawing \(R_{\rm pair} = \num{1000}\) uniformly random subsets \(\mathcal{T}'_{r}\) of \(|\mathcal{T}|\) pairs from the full set of \(P = N(N-1)/2\) pairs and recomputing \eqref{eq:methods_tissue_class} on each, the class-restricted trace is referred to the random-subset distribution,
\begin{equation}
    p^{\rm pair}_{\mathcal{T}} = \frac{\#\{\, r : \rhocoph_{\mathcal{T}'_{r}} \ge \rhocoph_{\mathcal{T}} \,\}}{R_{\rm pair}},
    \label{eq:methods_tissue_pairnull}
\end{equation}
so that the class \emph{concentrates} the trace when \(p^{\rm pair}_{\mathcal{T}} < \num{0.05}\) and is \emph{depleted} of it when \(p^{\rm pair}_{\mathcal{T}} > \num{0.95}\); matching \(|\mathcal{T}|\) makes the control immune to the noisier trace estimate of a small class. The second question is the mandatory gate: \(\rhocoph_{\mathcal{T}}\) is recomputed on each strength-preserving matched-strength surrogate (\(4\)-cycle \(\pm\delta\) rewiring, \(R\) surrogates; Section~\ref{sssec:methods_compare_stats}) restricted to the same class, and gated by the one-sided cohort Wilcoxon against the surrogate median, identically to the trace. Because the pair-count control preserves the connectivity that hyperconnected tissue inherits, it alone cannot separate a tissue effect from node strength; a class is read as carrying the trace only when it both concentrates by \eqref{eq:methods_tissue_pairnull} and clears the matched-strength gate.

Every class is evaluated per patient and summarized across the cohort: the pair-count control compares the median class trace across implanted patients against the same patient-median taken over each random-subset draw, and the matched-strength gate is the one-sided paired Wilcoxon of the per-patient class trace against its surrogate median, paired with the leave-one-patient-out diagnostic of Section~\ref{sssec:methods_compare_stats}. Because the \gls{imcoh} substrate is immune to zero-lag volume conduction, a tissue class that carries the trace is a connectivity grouping and not physical proximity, and no spatial null is applied.



\subsection{Temporal reinstatement of the trace}
\label{ssec:methods_reinstatement}

The primary trace treats post-task rest as one averaged network; resolving it in time asks whether the resting brain settles into the task configuration and dwells there or only passes through it in brief, replay-like excursions. Each resting recording (\acrshort{rspre}, \acrshort{rspost}) is cut into non-overlapping \qty{20}{\second} windows --- the shortest window at which the \(\alpha\) and \(\beta\) cophenetic image is itself reliable, below \(\sim\qty{10}{\second}\) the per-window estimate degrades and only \(\gamma\) supports faster windows --- and each window \(w\) is carried through the full pipeline of Section~\ref{ssec:methods_lrg} on its own \(\abs{\ImCoh}\) estimate, the windowed recipe reproducing the full-recording estimator in the full-window limit, to yield a per-window cophenetic image \(\Dcoph^{w}\). The window's task-likeness is the difference of two Spearman correlations of upper-triangular cophenetic distances,
\begin{equation}
    g_{w} = \rho_{S}\!\left(\Dcoph^{w}, \Dcoph^{\txtacr{taskt}}\right) - \rho_{S}\!\left(\Dcoph^{w}, \Dcoph^{\txtacr{rspre}}\right),
    \label{eq:methods_reinstatement_g}
\end{equation}
positive when the window's hierarchy resembles the task state more than the pre-task baseline, and each phase is summarized by its window mean \(A_{1}^{\Phase} = \operatorname{mean}_{w \in \Phase} g_{w}\). The cohort statistic is the one-sided paired Wilcoxon of the per-patient occupancy shift \(A_{1}^{\txtacr{rspost}} - A_{1}^{\txtacr{rspre}}\) under \(H_{1}\): post-task rest more task-like than pre-task rest, the per-patient value averaged over the two static-trace bands \(\alpha\) and \(\beta\) (the shift holds in each alone) and paired with the leave-one-patient-out diagnostic of Section~\ref{sssec:methods_compare_stats}. Because \(g_{w}\) is the same cophenetic quantity read at finer temporal resolution rather than a new connectivity measure, the reinstatement inherits the matched-strength gate of the static trace; a dedicated window-level surrogate is an optional supplementary control.

The occupancy shift is checked for scale- and representation-invariance and against the transient-replay alternative. It stays positive across the full diffusion-scale sweep of the propagator (every \(\tau\) from \(\tau_{\min}\) toward the Fiedler scale) and across five Laplacian-propagator read-outs of the same windows --- cophenetic, directional (magnetic-Laplacian), raw propagator, leading-subspace (Grassmann), and symmetric-normalized --- each positive in \(8\)--\(10\) of the \(10\) patients (\(p \le 0.007\)). The competing picture, that the task state re-enters rest as discrete, recurring bursts, is a cohort negative: window-to-window template burstiness (lag-one autocorrelation against a time-shuffled null), recurring-state discovery (\(k\)-means on the window cophenetic images, occupancy lift against the same null), and ordered sequence are each evaluated across every band, every diffusion scale, both task targets (\acrshort{taskt} and \acrshort{taskl}), the per-pair and \gls{ofc}-subtree resolutions, and every resolvable timescale from \qty{20}{\second} down to the \(\sim\qty{0.2}{\second}\) high-\(\gamma\) coherence floor, each against a time-shuffled and node-permuted placebo null. A single event at the ripple scale (\(\sim\qty{100}{\milli\second}\)) falls below the coherence floor of any connectivity estimator, so the transient reading is excluded only down to \(\sim\qty{0.2}{\second}\); and whether \acrshort{rspost} is drawn \emph{more tightly} toward the task state than \acrshort{rspre} is toward its own baseline is only a weak trend (paired \(p = 0.08\)), so we report a maintained proximity, not a sharpened attractor.

\subsection{Decomposing the trace into encoding and inference components}
\label{ssec:methods_arc}

The four-phase design --- \acrshort{rspre}, \acrshort{taskl}, \acrshort{taskt}, \acrshort{rspost} --- separates two processes the task combines. The patient first \emph{learns} a set of premise pairs by being shown them (\acrshort{taskl}), and then in the \emph{test} phase reasons about novel pairs that were never shown (\acrshort{taskt}). The primary trace of Section~\ref{ssec:methods_compare} is read from the test phase alone and therefore cannot distinguish what the network retains from \emph{seeing} the premises from what it retains of the relations it had to \emph{infer}; because the learning phase supplies a seen-only reference that requires no inference, the persistent reorganization splits into an encoding component and an inference-specific component. We reuse the split-baseline per-pair construction of Section~\ref{sssec:methods_compare_perpair} on the cophenetic communication distance \(O = \Dcoph\), writing \(\Dcoph^{\Phase}\) for the cophenetic image of phase \(\Phase\) and splitting \acrshort{rspre} into the same two independent halves \(A\) and \(B\). Alongside the task-side shift \(\Delta_{\rm task}^{A} = \Dcoph^{\txtacr{taskt}} - \Dcoph^{\txtacr{rspre}_{A}}\) that drives the primary trace, the design supplies an encoding shift, an inference-specific shift, and the persistence shift,
\begin{equation}
    \Delta_{\rm enc}^{A} = \Dcoph^{\txtacr{taskl}} - \Dcoph^{\txtacr{rspre}_{A}}, \qquad
    \Delta_{\rm inf} = \Dcoph^{\txtacr{taskt}} - \Dcoph^{\txtacr{taskl}}, \qquad
    \Delta_{\rm rest}^{B} = \Dcoph^{\txtacr{rspost}} - \Dcoph^{\txtacr{rspre}_{B}},
    \label{eq:methods_arc_shifts}
\end{equation}
where the encoding shift is baselined on one \acrshort{rspre} half and the persistence shift on the other, exactly as the task and rest sides of the primary trace, while the inference-specific shift is referenced to the \emph{learning} phase and so carries no \acrshort{rspre} half at all. The encoding shift measures how the communication hierarchy is displaced while the premises are shown; the inference-specific shift retains only the further displacement added during reasoning, with the encoding displacement already differenced out.

Each component is read through the arm-symmetric estimator of Section~\ref{sssec:methods_compare_perpair}: the two equally-valid assignments of the \acrshort{rspre} halves to the task and rest sides are averaged, so that the estimate is invariant to the arbitrary \(A\!\leftrightarrow\!B\) labelling, and the single-arm value is retained only as a supplement. The encoding trace is the symmetrized correlation of the encoding shift with persistence, in the rank form and sign orientation of the primary trace,
\begin{equation}
    T_{\rm enc} = \tfrac{1}{2}\Big[\rho_{S}\!\left(\Delta_{\rm enc}^{A},\, \Delta_{\rm rest}^{B}\right) + \big(A\!\leftrightarrow\!B\big)\Big],
    \label{eq:methods_arc_Tenc}
\end{equation}
where \((A\!\leftrightarrow\!B)\) denotes the mirror arm with the two \acrshort{rspre} halves exchanged; this is the cophenetic trace of Section~\ref{sssec:methods_compare_ctm} (equation~\eqref{eq:methods_rho_coph}) with the learning phase in place of the test phase, and asks whether the change induced by seeing the premises persists into post-task rest. The inference-specific trace asks whether the change added during reasoning persists once encoding is held constant, and removes the encoding contribution a second time as a first-order partial Spearman correlation, again arm-symmetrized,
\begin{equation}
    T_{\rm inf} = \tfrac{1}{2}\Big[\rho_{S}\!\left(\Delta_{\rm inf},\, \Delta_{\rm rest}^{B} \,\middle|\, \Delta_{\rm enc}^{A}\right) + \big(A\!\leftrightarrow\!B\big)\Big],
    \label{eq:methods_arc_Tinf}
\end{equation}
in which the first-order partial Spearman correlation of two shifts \(x,y\) controlling a third \(z\) is
\begin{equation}
    \rho_{S}\!\left(x, y \,\middle|\, z\right) = \frac{\rho_{S}(x, y) - \rho_{S}(x, z)\,\rho_{S}(y, z)}{\sqrt{\left[\,1 - \rho_{S}(x, z)^{2}\,\right]\left[\,1 - \rho_{S}(y, z)^{2}\,\right]}},
    \label{eq:methods_arc_partial}
\end{equation}
each \(\rho_{S}\) a Spearman coefficient across all \(N(N-1)/2\) contact pairs. Encoding is thus removed twice --- once in the reference of \(\Delta_{\rm inf}\) to the learning phase, and again by partialling \(\Delta_{\rm enc}\) out of the correlation --- which is what makes \(T_{\rm inf}\) inference-specific: it isolates persistence tied to the relations the patient inferred and not to the premises the patient was shown.

The two component traces symmetrize on different axes. The encoding trace \(T_{\rm enc}\), whose two shifts each carry an \acrshort{rspre} half, symmetrizes on both, exactly as the primary trace does. The inference-specific shift \(\Delta_{\rm inf}\), by contrast, is referenced to no \acrshort{rspre} half and is therefore invariant under the \(A\!\leftrightarrow\!B\) exchange, so the symmetrization of \(T_{\rm inf}\) acts only on the two shifts that do carry a half --- the persistence shift and the encoding shift that conditions it --- a narrower average than for \(T_{\rm enc}\). In both cases the arm-symmetric estimator is computed identically on the observation and on every matched-strength surrogate, and each component trace is gated per band by the matched-strength surrogate null and the one-sided cohort Wilcoxon of Section~\ref{sssec:methods_compare_stats}, identically to the primary trace, with positive values in the trace direction.

\paragraph{Scale dependence.}
Both component traces are read at the single working scale \(\tau_{\min} = 1/\lambda_{\max}\), as throughout (Section~\ref{ssec:methods_lrg}). Larger \(\tau\) weights longer diffusion walks --- the multi-step paths along which a transitive order chains --- which might invite re-reading the decomposition at coarser scales; but on the dense, continuous-spectrum \(\abs{\ImCoh}\) substrate such a \(\tau\)-sweep is degenerate. The propagator carries no spectral gap to single out an interior scale, and as \(\tau\) grows it relaxes monotonically toward its uniform mode until no hierarchy remains to read. The multiscale content of the decomposition is instead carried by the nested merge-height continuum of the cophenetic image \(\Dcoph\), which already folds every coarsening scale of \(\propag(\tau_{\min})\) into one per-pair hierarchy at the working scale (Section~\ref{sssec:methods_compare_ctm}); we therefore read \(T_{\rm enc}\) and \(T_{\rm inf}\) at \(\tau_{\min}\) alone.

\paragraph{Recording-length control.}
The test recording is longer than the learning recording, and the inference-specific shift \(\Delta_{\rm inf} = \Dcoph^{\txtacr{taskt}} - \Dcoph^{\txtacr{taskl}}\) inherits that asymmetry, so a component that merely grew with recording length would confound it. We test the dependence directly, regressing the per-patient arm-symmetric \(T_{\rm inf}\) on the per-patient test-to-learn duration ratio across the cohort by Spearman correlation. The encoding trace, built from the shorter learning recording against baseline alone, carries none of the test-versus-learning length difference and is duration-immune by construction.

\subsection{Epileptogenic-zone localization from the diffusion propagator}
\label{ssec:methods_epi}

\subsubsection{Relational seizure-onset affinity from the within-phase propagator}
\label{sssec:methods_epi_affinity}
The cross-phase trace of Section~\ref{ssec:methods_compare} measures how far the task's reorganization of the diffusion hierarchy persists into rest; that hierarchy is generated by a propagator on the connectivity graph, and the same propagator, read directly within a single recording, supports a second read-out that localizes epileptogenic tissue. On each patient's resting (post-task) per-band graph \(\Adjacency = \abs{\ImCoh}\) --- symmetric, non-negative, zero-diagonal, one node per \gls{seeg} contact --- we form the combinatorial Laplacian \(\lapl = \degm - \Adjacency\), with \(\degm = \diag(s)\) and node strength \(s_i = \sum_{j}\Adjacency_{ij}\), and its heat-kernel propagator
\begin{equation}
    \propag(\tau) = e^{-\tau\lapl} = U\,\diag\!\big(e^{-\tau\lambda_{k}}\big)\,U^{\transp}, \qquad \dm(\tau) = \propag(\tau)/\partf, \quad \partf = \Tr\,e^{-\tau\lapl},
    \label{eq:methods_epi_kernel}
\end{equation}
with \((\lambda_{k},U)\) the eigendecomposition of \(\lapl\) (Section~\ref{ssec:methods_lrg}). Whereas the trace reads the normalized operator \(\dm(\tau)\) at the single fine scale \(\tau_{\min}=1/\lambda_{\max}\) and compares phases, the marker reads a single entry \(\big[\propag(\tau)\big]_{ij}\) --- the heat that reaches contact \(i\) from contact \(j\) in diffusion time \(\tau\) --- within one phase, so none of the communication-distance, dendrogram, cophenetic, or Spearman machinery of Section~\ref{ssec:methods_compare} enters here, and because the global partition function \(\partf\) is a scalar that leaves the ranking of contacts unchanged, the affinity is read from the unnormalized kernel. Scale is swept over a per-patient logarithmic grid of six diffusion times,
\begin{equation}
    \tau_{\ell} = \frac{10^{\,\ell/5}}{\lambda_{\max}}, \qquad \ell = 0,1,\dots,5,
    \label{eq:methods_epi_tau}
\end{equation}
so that \(\tau_{0}=1/\lambda_{\max}=\tau_{\min}\) is the fine scale of the trace and \(\tau_{5}=10/\lambda_{\max}\) the coarse scale at which the seizure marker is read. Reading the marker at this coarse scale is dictated by the object it measures and licensed by its within-phase construction. Seizure-onset tissue is distinguished not by a single strong edge but by how broadly its influence diffuses through the network --- a mesoscale reachability resolved only once the propagator has integrated paths well beyond the immediate neighbourhood, to which the fine scale \(\tau_{\min}\) that anchors the cross-phase trace is blind. The trace is confined to \(\tau_{\min}\) because coarsening a cross-phase comparison drives its diffusion geometry toward the trivial mode, where a phase-invariant baseline dominates every phase alike; a within-phase affinity ranking has no cross-phase contrast to lose and may therefore be read at the coarse end of the resolution window \([\tau_{\min}, \tau^{\ast}]\) of Section~\ref{ssec:methods_lrg}. Accordingly \(\tau_{5}=10/\lambda_{\max}\) sits near the entropic-susceptibility peak \(\tau^{\ast}\), the coarsest scale at which inter-contact structure survives before the propagator collapses onto its uniform mode. Let \(E\subset V\) be the set of seizure-onset contacts marked by the clinician on the implantation record. \(E\) is the ground-truth label set: it defines the positives against which the marker is scored and is never handed to the marker whole. The marker is instead \emph{seeded} --- given a set \(S\subseteq E\) of \emph{known} onset contacts, the affinity of a contact \(i\) to a target set \(T\) is the mean kernel weight reaching it from that set,
\begin{equation}
    \operatorname{aff}(i,T) = \frac{1}{\abs{T\setminus i}}\sum_{s\,\in\,T\setminus i}\big[\propag(\tau)\big]_{is}, \qquad \operatorname{aff}(i)\equiv\operatorname{aff}(i,S),
    \label{eq:methods_epi_aff}
\end{equation}
and the seizure read-out is the seed affinity \(\operatorname{aff}(i,S)\). The marker ranks the non-seed contacts by this score and aims to place the \emph{unseeded} onset contacts \(E\setminus S\) above the healthy contacts \(V\setminus E\), so that the seizure zone is recovered from a handful of known seeds rather than assumed; which contacts constitute the seed \(S\) is fixed by the evaluation (Section~\ref{sssec:methods_epi_gate}) --- every \emph{other} onset contact when the question is whether each onset site is recovered by its peers, a few known seeds for the deployable detector, or the onset contacts of \emph{other} electrodes for distant discovery. We read the propagator \emph{relationally}: a node-intrinsic reading --- each contact summarized by its own diffusion through the return probability \(\big[\propag(\tau)\big]_{ii}\), the diffusion-row entropy, or the commute distance --- is defined only as a comparison, and the marker is the relational seed affinity rather than any single-node quantity. Because seizure tissue tends to be hyperconnected, node strength alone carries part of this signal; to keep the read-out orthogonal to hubness we residualize the affinity on strength, fitting \(\operatorname{aff}(i)\approx a\,s_{i}+b\) by least squares across contacts and retaining the residual
\begin{equation}
    \widetilde{\operatorname{aff}}(i) = \operatorname{aff}(i) - \big(a\,s_{i}+b\big).
    \label{eq:methods_epi_resid}
\end{equation}
Since the \gls{imcoh} substrate is immune to zero-lag volume conduction, this affinity cannot be a restatement of physical proximity, and a strength-matched null is the only control the connectivity read-out requires.

\subsubsection{Scoring and the strength-matched gate}
\label{sssec:methods_epi_gate}
Each contact score is turned into a marker by rank concordance with the labels: the \gls{auc} is the Mann--Whitney probability that a seizure contact outranks a healthy one,
\begin{equation}
    \mathrm{AUC} = \frac{\#\{\,\text{seizure}>\text{healthy}\,\} + \tfrac12\,\#\{\,\text{ties}\,\}}{n_{\text{seizure}}\;n_{\text{healthy}}}.
    \label{eq:methods_epi_auc}
\end{equation}
The primary evaluation is within patient and uses every marked contact: holding out one onset contact \(i\) at a time, we take the seed set to be its peers, \(S = E\setminus\{i\}\), and rank the held-out contact against every healthy contact by the strength-residualized score \(\widetilde{\operatorname{aff}}\) of \eqref{eq:methods_epi_resid}, so that no contact ever contributes to its own score; a patient is admitted only when at least four seizure and five healthy contacts are present. Significance is assessed against a strength-matched label null: \num{300} surrogate label sets, each matched to \(E\) in cardinality and in node-strength quartile, are passed through the identical leave-one-out marker to yield a per-patient upper-tail \(p\) and a cohort empirical \(p\) from the across-patient median score, and the cohort effect is a one-sided Wilcoxon signed-rank test of the per-patient \(\mathrm{AUC}-\tfrac12\). No spatial or proximity null is applied: it would be redundant on a substrate immune to zero-lag volume conduction, where matched strength is the mandatory control. As a conservative test of \emph{distant} discovery, we additionally evaluate off-shaft, partitioning contacts by electrode: for each electrode that carries seizure contacts we seed only on the seizure contacts of \emph{other} electrodes and score only contacts lying on electrodes that carry no seed, so that adjacency to a known contact cannot contribute and the read-out is tested where proximity is useless by construction.

\subsubsection{A fused, coordinate-free seizure-zone probability}
\label{sssec:methods_epi_detector}
To combine scales and rhythms into a single deployable read-out, we summarize each contact, relative to \(k=\num{3}\) known seizure seeds \(S\), by a vector of diffusion features: the heat-kernel affinity of \eqref{eq:methods_epi_aff} at fine, intermediate, and coarse \(\tau\), together with a small family of standard graph-diffusion kernels each read as a seed affinity --- personalized PageRank \((1-\eta)\big(I-\eta\,\degm^{-1}\Adjacency\big)^{-1}\) with damping \(\eta\), the Katz resolvent \((I-\kappa\Adjacency)^{-1}\) at sub-critical attenuation \(\kappa\), communicability \(e^{\Adjacency}\), the negative commute-time (effective resistance) \(-R\) with \(R_{ij}=\lapl^{+}_{ii}+\lapl^{+}_{jj}-2\lapl^{+}_{ij}\), and a diffusion distance between heat profiles (the full kernel set is listed in the Supplement) --- and the seed-relative segregation \(\operatorname{aff}(i,S)-\operatorname{aff}(i,V\setminus S)\). No contact coordinates enter the feature set. The per-band features are concatenated across all six frequency bands, standardized within patient, and passed to an \(L_{2}\)-regularized logistic model trained leave-one-patient-out, producing a calibrated seizure-onset probability
\begin{equation}
    P(\mathrm{SOZ}\mid i) = \sigma\!\Big(b_{0} + \textstyle\sum_{f} b_{f}\,z_{f}(i)\Big),
    \label{eq:methods_epi_detector}
\end{equation}
with \(\sigma\) the logistic function and \(z_{f}\) the standardized features. The detector is scored by cross-patient \gls{auc}, by precision among its five highest-probability contacts and recall within its ten highest, and by the Brier score, all referred to the same strength-matched label shuffle as an empirical floor; we report the interpretable logistic model in preference to a higher-capacity nonlinear classifier, which is prone to overfitting at this cohort size.

\subsubsection{Seed-free features, the two-population switch, and validation}
\label{sssec:methods_epi_seedfree}
Two further reductions complete the analysis. First, a seed-free variant removes the dependence on known contacts, replacing the seed affinity with each contact's intrinsic diffusion signature --- the self-return \(\big[\propag(\tau)\big]_{ii}\), the localization of the slow Laplacian modes at the node (the inverse participation ratio of the low-\(\lambda\) eigenvectors), and its cophenetic persistence --- alongside a purely geometric baseline, the normalized along-shaft electrode depth; these are fused across bands and transferred leave-one-patient-out exactly as in \eqref{eq:methods_epi_detector}. Second, because a seizure network can present either as a co-diffusing community, in which the relational affinity is informative, or as a single high-strength hub, in which it is not, we select per patient between the two read-outs with a hard switch on the mean strength percentile \(r\) of the seeds,
\begin{equation}
    s_{\mathrm{switch}}(i) = \begin{cases} s_{i} & r \ge 0.70 \quad(\text{hub regime}),\\[2pt] \widetilde{\operatorname{aff}}(i) & r < 0.70 \quad(\text{community regime}), \end{cases}
    \label{eq:methods_epi_switch}
\end{equation}
rather than a continuous blend, which would average a working detector against a known failure mode. Performance is accordingly reported per patient and never pooled into a single cohort number. All read-outs are validated against the clinician-marked seizure-onset labels rather than against resection margins or surgical outcome, so the read-out is characterized as a connectivity marker that co-locates with the clinical seizure-onset zone, its precision interpreted as a triage shortlist over the implanted contacts.